\documentclass[12pt,a4paper]{article}

% ----- Encoding and language -----------------------------------------
\usepackage[utf8]{inputenc}
\usepackage[T1]{fontenc}
\usepackage{lmodern}
\usepackage[spanish,es-tabla,es-nodecimaldot, es-noquoting]{babel}
\usepackage{microtype}

% ----- Page layout --------------------------------------------------- 
\usepackage[a4paper,left=2.6cm,right=2.6cm,top=2.8cm,bottom=2.6cm,headheight=15pt]{geometry}
\linespread{1.15}
\setlength{\parskip}{4pt}
\setlength{\parindent}{0pt}

% ----- Colors --------------------------------------------------------
\usepackage[table]{xcolor}
\definecolor{darkblue}{HTML}{0B2545}
\definecolor{mediumblue}{HTML}{1D4E89}
\definecolor{lightblue}{HTML}{DCE7F3}
\definecolor{borderblue}{HTML}{B8C7DD}
\definecolor{lightgray}{HTML}{F5F6F8}

% ----- Packages ------------------------------------------------------
\usepackage{graphicx,float,booktabs,tabularx,array,multirow,longtable}
\usepackage{enumitem,amsmath,amssymb,caption,xparse,url}
\usepackage{siunitx}
\sisetup{output-decimal-marker={,}, inter-unit-product=\ensuremath{{}\cdot{}}}
\usepackage{tikz}
\usetikzlibrary{arrows.meta,calc,positioning}
\usepackage{pgfplots}
\pgfplotsset{compat=1.16}
\usepackage[most]{tcolorbox}
\usepackage{fancyhdr,titlesec}

% ----- Header and footer ---------------------------------------------
\pagestyle{fancy}
\fancyhf{}
\fancyfoot[L]{\sffamily\footnotesize\bfseries\color{darkblue} Facultad de
Ingeniería Industrial}
\fancyfoot[R]{\sffamily\footnotesize\bfseries\color{darkblue}Pag \thepage}
\renewcommand{\headrulewidth}{0pt}
\renewcommand{\footrulewidth}{0.8pt}
\renewcommand{\footrule}{{\color{darkblue}\hrule width\headwidth
height\footrulewidth}}

% ----- Section titles -----------------------------------------------
\titleformat{\section}
  {\sffamily\bfseries\Large\color{darkblue}}{\thesection.}{0.6em}{}
  % [{\vspace{-0.6ex}\color{mediumblue}\titlerule[1.3pt]}]
\titleformat{\subsection}
  {\sffamily\bfseries\large\color{mediumblue}}{\thesubsection}{0.6em}{}
\titleformat{\subsubsection}
  {\sffamily\bfseries\normalsize\color{darkblue}}{\thesubsubsection}{0.6em}{}
\titlespacing*{\section}{0pt}{2.4ex plus 1ex minus .2ex}{1.6ex}
\titlespacing*{\subsection}{0pt}{2.0ex plus .8ex minus .2ex}{1.0ex}
\titlespacing*{\subsubsection}{0pt}{1.6ex plus .5ex}{0.6ex}
\addto\captionsspanish{\renewcommand{\contentsname}{CONTENIDO}}

% ----- Captions and tables ------------------------------------------
\captionsetup{font=small,labelfont={bf,color=darkblue},labelsep=period,
              textfont={color=black!80},skip=6pt}
\renewcommand{\arraystretch}{1.22}
\setlist{itemsep=2pt,topsep=3pt,parsep=0pt}

% Table header cell (white text on dark blue)
\newcommand{\encab}[1]{\textcolor{white}{\sffamily\bfseries #1}}

% Highlighting data that the student must confirm
\newcommand{\verificar}[1]{\colorbox{yellow!45}{#1}}

% ----- Text boxes ----------------------------------------------------
\newtcolorbox{nota}[1][Nota]{colback=lightblue,colframe=mediumblue,coltitle=white,
  fonttitle=\sffamily\bfseries\small,title=#1,arc=2pt,boxrule=0.6pt,
  left=7pt,right=7pt,top=4pt,bottom=4pt,breakable}
\newtcolorbox{enunciado}{colback=lightblue,colframe=darkblue,arc=2pt,boxrule=0pt,
  leftrule=3pt,left=8pt,right=8pt,top=4pt,bottom=4pt,breakable,fontupper=\sffamily\small}

% ----- Image with placeholder box if file does not exist -------------
% Uso: \imagen[width][box height]{file}{description}
\NewDocumentCommand{\imagen}{O{0.6} O{5cm} m m}{%
  \IfFileExists{#3}{\includegraphics[width=#1\linewidth]{#3}}{%
    {\setlength{\fboxsep}{0pt}\fcolorbox{borderblue}{lightgray}{%
      \parbox[c][#2][c]{#1\linewidth}{\centering\sffamily\small\color{darkblue}%
      \textbf{[ Insertar imagen
]}\\[3pt]\texttt{\detokenize{#3}}\\[3pt]\footnotesize #4}}}}}

% ----- Links ---------------------------------------------------------
\usepackage[colorlinks=true,breaklinks=true,linkcolor=darkblue,urlcolor=mediumblue,citecolor=darkblue,
  pdftitle={Laboratorio N.º 2: Metrología},
  pdfauthor={Ahon Saenz, Maria Julia Armida; Rey de Castro Rupay, Luis
Alejandro},
  pdfsubject={Procesos de Manufactura — UNMSM}]{hyperref}
\urlstyle{same}
\makeatletter\g@addto@macro\UrlBreaks{\do\-}\makeatother
\Urlmuskip=0mu plus 2mu\relax
\setlength{\emergencystretch}{2.5em}
\newcommand{\refbib}[1]{[\ref{#1}]}

% =====================================================================
\begin{document}
% =====================================================================

% ===== COVER PAGE ====================================================
\begin{titlepage}
	\thispagestyle{empty}
	\begin{tikzpicture}[remember picture,overlay]
		\fill[darkblue] (current page.north west) rectangle ($(current page.north
			east)+(0,-4.4cm)$);
		\fill[mediumblue] ($(current page.north west)+(0,-4.4cm)$) rectangle
		($(current page.north east)+(0,-4.65cm)$);
		\fill[darkblue] (current page.south west) rectangle ($(current page.south
			east)+(0,1.1cm)$);
		\node[text=white,align=center,text width=16cm,font=\sffamily] at ($(current
			page.north)+(0,-2.25cm)$) {%
			{\Large\bfseries UNIVERSIDAD NACIONAL MAYOR DE SAN MARCOS}\\[5pt]
			{\small\itshape Universidad del Perú, Decana de América}\\[10pt]
			{\normalsize\bfseries ESCUELA PROFESIONAL DE INGENIERÍA INDUSTRIAL}};
	\end{tikzpicture}

	\vspace*{2cm}

	\begin{center}
		\imagen[0.22][4.3cm]{img/unmsm-logo.png}{}

		\vspace{0.8cm}
		{\sffamily\bfseries\color{darkblue}\fontsize{22}{26}\selectfont
			Laboratorio N\textdegree 2: Metrología\par}

		\vspace{0.4cm}

		\begin{tcolorbox}[colback=white,colframe=white,arc=3pt,width=0.9\linewidth,
				box align=center,left=12pt,right=12pt,top=8pt,bottom=8pt]
			\centering\sffamily
			\begin{tabular}{@{}>{\bfseries\color{darkblue}}r@{\hspace{1em}}l@{}}
				Asignatura: & Procesos de Manufactura (L1-3) \\[5pt]
				Docente:    & Mg. Ing. Víctor Rosales Urbano \\
			\end{tabular}
		\end{tcolorbox}

		\vspace{0.4cm}
		{\sffamily\bfseries\color{darkblue}\large Integrantes\par}
		\vspace{5pt}
		\sffamily\begin{tabular}{@{}l@{\hspace{2.5em}}r@{}}
			Ahon Saenz, Maria Julia Armida      & (23170001) \\
			Rey de Castro Rupay, Luis Alejandro & (22170163) \\
		\end{tabular}

		\vspace{1cm}

		\imagen[0.6][5.2cm]{img/portada.png}{Imagen de portada}
	\end{center}

	\vfill
	\begin{center}
		{\sffamily\bfseries\color{darkblue}\large Lima, 2026\par}
	\end{center}
\end{titlepage}
\setcounter{page}{2}

% ===== CONTENIDO =====================================================
\tableofcontents
\newpage

% ===== 1. PRESENTACIÓN ===============================================
\section{PRESENTACIÓN}

La metrología es la ciencia de la medición y de sus aplicaciones. En manufactura
es la condición previa de casi todo lo demás: no se puede controlar una
tolerancia, garantizar que dos piezas ensamblen entre sí ni decidir si un lote
se acepta o se rechaza sin antes medir bien. En Ingeniería Industrial esto es
especialmente importante, porque las decisiones sobre calidad, costos y
productividad se apoyan en datos que solo son útiles si provienen de
mediciones confiables.\\[-4pt]

En este segundo laboratorio de Procesos de Manufactura medimos una misma pieza
de acero con tres instrumentos de distinta resolución: una regla (escuadra), un
vernier (pie de rey) y un micrómetro de exteriores. Lo valioso de la práctica es
que la pieza es siempre la misma y solo cambia la herramienta, de modo que las
diferencias entre lecturas pueden atribuirse al instrumento y al método de
medición, y no a la pieza. Esa comparación directa nos permitió ver, con
nuestros propios números, por qué existen instrumentos distintos y cuándo
conviene emplear cada uno.\\[-4pt]

\begin{figure}[H]
	\centering
	\imagen[0.45][4.2cm]{img/presentacion.jpeg}{}
\end{figure}

Además, el laboratorio nos obligó a trabajar con dos sistemas de unidades
(milímetros y pulgadas), algo habitual en la industria, donde planos,
herramientas y catálogos conviven en ambos sistemas. Un error de conversión
puede parecer trivial en el papel, pero, como se menciona en la sección de
accidentes relacionados, ha tenido consecuencias muy costosas en proyectos
reales.

% ===== 2. MARCO TEÓRICO ==============================================
\section{MARCO TEÓRICO}

\subsection{Sistema Internacional de Unidades}

El \textbf{Sistema Internacional de Unidades (SI)} es el sistema de medida
coherente de uso mundial. Su mantenimiento está a cargo de la Oficina
Internacional de Pesas y Medidas (BIPM), organismo creado a partir de la
Convención del Metro de 1875.\\[-4pt]

Desde el 20 de mayo de 2019 el SI se define fijando el valor numérico exacto de
siete constantes de la naturaleza (Tabla~\ref{tab:constantes}); con ello,
ninguna unidad depende ya de un objeto material, como ocurría con el antiguo
prototipo internacional del kilogramo. En el Perú, el Sistema Legal de Unidades
de Medida (SLUMP), establecido por la Ley N\textdegree 23560, adopta el SI junto
con sus múltiplos y submúltiplos decimales.\\[-4pt]

\begin{table}[!htb]
	\centering
	\caption{Constantes definitorias del SI}\label{tab:constantes}
	\small
	\rowcolors{2}{white}{lightblue}
	\begin{tabularx}{\linewidth}{>{\raggedright\arraybackslash}X c c c}
		\rowcolor{darkblue}
		\encab{Constante}                                    & \encab{Símbolo}
		                                                     & \encab{Valor exacto}
		                                                     & \encab{Unidad que
			                                                       fija}
		\\
		Frecuencia de la transición hiperfina del $^{133}$Cs &
		$\Delta\nu_{\mathrm{Cs}}$                            & $9\,192\,631\,770$ Hz
		                                                     & segundo
		\\
		Velocidad de la luz en el vacío                      & $c$
		                                                     & $299\,792\,458$ m/s
		                                                     & metro
		\\
		Constante de Planck                                  & $h$
		                                                     &
		$6{,}626\,070\,15\times10^{-34}$ J$\cdot$s           & kilogramo
		\\
		Carga elemental                                      & $e$
		                                                     &
		$1{,}602\,176\,634\times10^{-19}$ C                  & amperio
		\\
		Constante de Boltzmann                               & $k$
		                                                     &
		$1{,}380\,649\times10^{-23}$ J/K                     & kelvin
		\\
		Constante de Avogadro                                & $N_A$
		                                                     &
		$6{,}022\,140\,76\times10^{23}$ mol$^{-1}$           & mol
		\\
		Eficacia luminosa a $540\times10^{12}$ Hz            & $K_{\mathrm{cd}}$
		                                                     & $683$ lm/W            &
		candela
		\\
	\end{tabularx}
\end{table}

\vspace{4pt}

Las unidades pueden expresarse con prefijos decimales, de los cuales los más
usados en metrología dimensional se resumen en la Tabla~\ref{tab:prefijos}. Cabe
añadir que, fuera del SI, la industria sigue utilizando la \textbf{pulgada},
cuya equivalencia con el sistema métrico es exactamente $1\ \text{in} =
	\SI{25.4}{\milli\metre}$.\\[-4pt]

\begin{table}[!htb]
	\centering
	\caption{Prefijos del SI frecuentes en metrología
		dimensional}\label{tab:prefijos}
	\small
	\rowcolors{2}{white}{lightblue}
	\begin{tabular}{l c c @{\hspace{2.5em}} l c c}
		\rowcolor{darkblue}
		\encab{Prefijo} & \encab{Símbolo} & \encab{Factor} & \encab{Prefijo} &
		\encab{Símbolo} & \encab{Factor}
		\\
		kilo            & k               & $10^{3}$       & centi           & c
		                & $10^{-2}$                                                \\
		mega            & M               & $10^{6}$       & mili            & m
		                & $10^{-3}$
		\\
		giga            & G               & $10^{9}$       & micro           &
		\textmu         & $10^{-6}$
		\\
		deci            & d               & $10^{-1}$      & nano            & n &
		$10^{-9}$                                                                  \\
	\end{tabular}
\end{table}

\subsection{Unidades fundamentales}

Las \textbf{unidades fundamentales o de base} son aquellas a partir de las
cuales se construyen todas las demás. El SI reconoce siete, cada una asociada a
una magnitud física independiente (Tabla~\ref{tab:base}).

En este laboratorio la magnitud de interés es la \textbf{longitud}, cuya unidad
es el metro. En la práctica de taller se trabaja con su submúltiplo, el
milímetro ($\SI{1}{\milli\metre} = 10^{-3}\ \text{m}$), y con el micrómetro o
micra ($\SI{1}{\micro\metre} = 10^{-6}\ \text{m}$) cuando se habla de
tolerancias finas y de rugosidad.

\begin{table}[!htb]
	\centering
	\caption{Magnitudes y unidades de base del SI}\label{tab:base}
	\small
	\rowcolors{2}{white}{lightblue}
	\begin{tabular}{l l c }
		\rowcolor{darkblue}
		\encab{Magnitud}          & \encab{Unidad} & \encab{Símbolo} \\
		Longitud                  & metro          & m               \\
		Masa                      & kilogramo      & kg              \\
		Tiempo                    & segundo        & s               \\
		Corriente eléctrica       & amperio        & A               \\
		Temperatura termodinámica & kelvin         & K               \\
		Cantidad de sustancia     & mol            & mol             \\
		Intensidad luminosa       & candela        & cd              \\
	\end{tabular}
\end{table}

\subsection{Unidades derivadas}

Las \textbf{unidades derivadas} se obtienen combinando unidades de base mediante
productos y cocientes, sin factores numéricos adicionales (por eso el SI es un
sistema \emph{coherente}). Algunas reciben un nombre especial en honor a un
científico. La Tabla~\ref{tab:derivadas} resume las que más aparecen en
manufactura.\\[-4pt]

\begin{table}[!htb]
	\centering
	\caption{Unidades derivadas del SI de uso frecuente en
		manufactura}\label{tab:derivadas}
	\small
	\rowcolors{2}{white}{lightblue}
	\begin{tabularx}{\linewidth}{l l c >{\raggedright\arraybackslash}X}
		\rowcolor{darkblue}
		\encab{Magnitud}        & \encab{Unidad}             & \encab{Símbolo} &
		\encab{Expresión}
		\\
		Área                    & metro cuadrado             & m$^2$           & m
		$\cdot$ m
		\\
		Volumen                 & metro cúbico               & m$^3$           & m
		$\cdot$ m $\cdot$ m
		\\
		Velocidad               & metro por segundo          & m/s             & m
		$\cdot$ s$^{-1}$
		\\
		Velocidad angular       & radián por segundo         & rad/s           &
		s$^{-1}$ (1 rad $=$ m/m)
		\\
		Frecuencia              & hercio                     & Hz              &
		s$^{-1}$
		\\
		Densidad                & kilogramo por metro cúbico & kg/m$^3$        & kg
		$\cdot$ m$^{-3}$
		\\
		Fuerza                  & newton                     & N               & kg
		$\cdot$ m $\cdot$ s$^{-2}$
		\\
		Presión                 & pascal                     & Pa              &
		N/m$^2$ $=$ kg $\cdot$ m$^{-1}$ $\cdot$ s$^{-2}$
		\\
		Energía, trabajo        & julio                      & J               & N
		$\cdot$ m
		\\
		Potencia                & vatio                      & W               & J/s
		\\
		Par (momento de fuerza) & newton metro               & N$\cdot$m       & N
		$\cdot$ m
		\\
		Tensión eléctrica       & voltio                     & V               & W/A
		\\
		Resistencia eléctrica   & ohmio                      & $\Omega$        & V/A
		\\
	\end{tabularx}
\end{table}

\vspace{4pt}

Aunque el par y la energía son dimensionalmente iguales (N$\cdot$m), el SI
reserva el julio para la energía y expresa el par como newton metro. Por otro
lado, magnitudes como el área y el volumen de una pieza se obtienen
\emph{indirectamente} a partir de longitudes medidas con vernier, micrómetro o
algún otro instrumento de medición de longitud.

\subsection{Instrumentos de medición}

Un \textbf{instrumento de medición} es un dispositivo utilizado para realizar
mediciones, solo o acompañado de dispositivos complementarios. Para comparar
instrumentos conviene manejar los siguientes conceptos:

\vspace{4pt}

\begin{itemize}
	\item \textbf{Resolución:} menor variación de la magnitud que produce una
	      variación perceptible en la indicación del instrumento.
	\item \textbf{Exactitud:} cercanía entre el valor medido y el valor de
	      referencia.
	\item \textbf{Precisión:} cercanía entre mediciones repetidas en las mismas
	      condiciones.
	\item \textbf{Error de medición:} diferencia entre el valor medido $x_m$ y un
	      valor de referencia $x_{ref}$. Se expresa como error absoluto y
	      relativo:
	      \[
		      E_a = x_m - x_{ref}, \qquad E_r = \frac{|E_a|}{x_{ref}}\times
		      100\,\text{\%}
	      \]
	\item \textbf{Calibración y trazabilidad:} comparación del instrumento con
	      patrones cuyo valor está vinculado, mediante una cadena ininterrumpida,
	      a los patrones nacionales e internacionales.
\end{itemize}

\vspace{4pt}

Los instrumentos de longitud utilizados en el laboratorio son de \emph{lectura
	directa sobre una escala}. Sus principios y resoluciones se resumen a
continuación.

\vspace{4pt}

\begin{description}[leftmargin=0pt,style=unboxed,labelindent=0pt]
	\item[Regla y escuadra graduadas.] Se lee la posición del extremo de la pieza
	      sobre una escala impresa o grabada. Con divisiones de
	      \SI{1}{\milli\metre}, la
	      resolución es de \SI{1}{\milli\metre} y la incertidumbre de lectura se
	      estima en la mitad de una división ($\pm\SI{0.5}{\milli\metre}$), sin
	      contar errores de paralaje ni de alineación.\\[-4pt]
	\item[Vernier o pie de rey.] Usa una segunda escala deslizante (el
	      \emph{nonio}) con $n$ divisiones que abarcan $n-1$ divisiones de la
	      escala principal, lo que permite leer fracciones de división. Si $d$ es
	      la menor división de la escala principal:
	      \[
		      r = \frac{d}{n} = \frac{\SI{1}{\milli\metre}}{20} =
		      \SI{0.05}{\milli\metre},
		      \qquad L = L_{\text{principal}} + k\cdot r
	      \]
	      donde $k$ es el número de la división del nonio que coincide con una
	      raya de la escala principal. Los vernier con escala en pulgadas de
	      $1/128$ in resuelven aproximadamente \SI{0.198}{\milli\metre}.\\[-4pt]
	\item[Micrómetro de exteriores.] Se basa en el avance de un tornillo de paso
	      conocido. Si el paso es $p = \SI{0.5}{\milli\metre}$ y el tambor tiene
	      $N = 50$
	      divisiones, cada división equivale a
	      \[
		      r = \frac{p}{N} = \frac{\SI{0.5}{\milli\metre}}{50} =
		      \SI{0.01}{\milli\metre}
	      \]
	      La lectura es la suma de la lectura del manguito (escala fija) y la del
	      tambor. El trinquete garantiza una fuerza de medición constante, y el
	      rango típico de cada micrómetro es de \SI{25}{\milli\metre}.
\end{description}

\begin{nota}[Criterio práctico de selección]
	La llamada \emph{regla de diez a uno} recomienda que la resolución del
	instrumento sea, como mucho, una décima parte de la tolerancia que se desea
	verificar. Así, para una tolerancia de $\pm\SI{0.1}{\milli\metre}$ (banda de
	\SI{0.2}{\milli\metre}) se requiere resolución de \SI{0.02}{\milli\metre} o
	mejor, lo que descarta la regla y obliga a usar un micrómetro.
\end{nota}

% ===== 3. MATERIALES Y EQUIPOS =======================================
\section{MATERIALES Y EQUIPOS}

\subsection{Regla}

Se empleó una \textbf{escuadra} de plástico transparente con escala graduada en
milímetros. Es el instrumento más simple y económico del laboratorio y sirve
para mediciones aproximadas: al no tener un tope de referencia ni mordazas, la
lectura depende de que el operador alinee bien el cero con el borde de la pieza
y observe la escala de frente para evitar el error de paralaje.

\begin{figure}[H]
	\centering
	\imagen[0.55][4.2cm]{img/regla.png}{}
	\caption{Regla (escuadra) utilizada en las mediciones.}
\end{figure}

\begin{table}[!htb]
	\centering
	\small
	\rowcolors{1}{white}{lightblue}
	\begin{tabularx}{\linewidth}{>{\bfseries}l X}
		Tipo                     & Escuadra plástica graduada
		\\
		Graduación / resolución  & \SI{1}{\milli\metre}
		\\
		Incertidumbre de lectura & $\pm\SI{0.5}{\milli\metre}$ (media división), más
		paralaje y alineación
		\\
		Rango                    & aprox. 0--200 mm
		\\
	\end{tabularx}
\end{table}

\subsection{Vernier o Pie de Rey}

El vernier utilizado es de marca Mitutoyo, con escala principal en milímetros y
en pulgadas. Consta de mordazas exteriores (para longitudes y espesores),
mordazas interiores (para diámetros de agujeros y ranuras), una
varilla de profundidad, la escala principal, el nonio y un tornillo de fijación
de lectura. Las lecturas registradas son múltiplos de \SI{0.05}{\milli\metre} en
la escala métrica y de $1/128$ in en la escala inglesa.

\begin{figure}[H]
	\centering
	\imagen[0.5][4.6cm]{img/vernier.jpeg}{}
	\caption{Vernier o pie de rey (Mitutoyo).}
\end{figure}

\begin{table}[!htb]
	\centering
	\small
	\rowcolors{1}{white}{lightblue}
	\begin{tabularx}{\linewidth}{>{\bfseries}l X}
		Marca      & Mitutoyo
		\\
		Resolución & \SI{0.05}{\milli\metre} (escala métrica) y $1/128$ in (escala
		en pulgadas)
		\\
		Rango      & aprox. 0--150 mm
		\\
		Aplicación & Exteriores, interiores (diámetros de agujeros) y profundidades
		\\
	\end{tabularx}
\end{table}

\subsection{Micrómetro}

Se utilizó un micrómetro de exteriores compuesto por arco, yunque, husillo,
manguito, tambor y trinquete. Su rango de medición es de \SI{0}{} a
\SI{25}{\milli\metre} y su resolución de \SI{0.01}{\milli\metre}, la mejor de
los tres instrumentos. Por la profundidad limitada de su arco y su rango máximo,
solo permite medir dimensiones accesibles desde el borde de la pieza.\\[-4pt]

\begin{figure}[H]
	\centering
	\imagen[0.55][4.6cm]{img/micrometro.jpeg}{Micrómetro de exteriores (0--25 mm)}
	\caption{Micrómetro de exteriores empleado en el laboratorio.}
\end{figure}

\begin{table}[!htb]
	\centering
	\small
	\rowcolors{1}{white}{lightblue}
	\begin{tabularx}{\linewidth}{>{\bfseries}l X}
		Tipo       & Micrómetro de exteriores con trinquete                     \\
		Rango      & \SIrange{0}{25}{\milli\metre}                              \\
		Resolución & \SI{0.01}{\milli\metre}                                    \\
		Limitación & Profundidad del arco: solo mide bordes accesibles; no mide
		diámetros de agujeros                                                   \\
	\end{tabularx}
\end{table}

\subsection{Pieza de acero}

La pieza medida es una placa de acero de acabado oscuro, con un contorno
escalonado, una ranura rectangular en uno de sus extremos y tres agujeros
pasantes. Para identificar sus dimensiones se asignó una letra a cada lado y a
cada diámetro (sección~\ref{sec:resultados}). Sus dimensiones generales son del
orden de $50\times40\times\SI{11.65}{\milli\metre}$.\\[-4pt]

\begin{figure}[H]
	\centering
	\imagen[0.6][5cm]{img/pieza.jpeg}{Pieza de acero con escalón, ranura y tres
		agujeros}
	\caption{Pieza de acero sometida a medición.}
\end{figure}

% ===== 4. PROCEDIMIENTO ==============================================
\section{PROCEDIMIENTO}

La secuencia de trabajo se resume en la Figura~\ref{fig:flujo} y se detalla a
continuación.\\[-4pt]

\begin{figure}[H]
	\centering
	\begin{tikzpicture}[node distance=9mm,
			caja/.style={rectangle,rounded
					corners=3pt,fill=darkblue,text=white,font=\sffamily\footnotesize\bfseries,
					minimum height=1.15cm,minimum width=2.55cm,align=center,inner sep=3pt},
			flecha/.style={-{Stealth[length=2.2mm]},thick,color=mediumblue}]
		\node[caja] (a) {Identificar\\lados y diámetros};
		\node[caja,right=of a] (b) {Medir con\\regla};
		\node[caja,right=of b] (c) {Medir con\\vernier (mm y in)};
		\node[caja,right=of c] (d) {Medir con\\micrómetro};
		\draw[flecha] (a)--(b); \draw[flecha] (b)--(c); \draw[flecha] (c)--(d);
	\end{tikzpicture}
	\caption{Secuencia del procedimiento experimental.}\label{fig:flujo}
\end{figure}

\begin{enumerate}[label=\textbf{\arabic*.}]
	\item Se asignó una letra a cada lado de la figura, así como a los diámetros
	      de los agujeros presentes, a modo de identificación
	      (Figura~\ref{fig:croquis}).
	\item Se midió la longitud de cada una de estas partes con una regla
	      (escuadra) y se registraron los resultados en una tabla.
	\item Se repitió el mismo procedimiento con un vernier, registrando las
	      medidas en milímetros y en pulgadas en dos columnas distintas.
	\item Finalmente se empleó el micrómetro, aunque solo fue posible registrar
	      dos mediciones debido a las limitaciones del instrumento (estructura
	      física y
	      rango máximo de \SI{25}{\milli\metre}).
\end{enumerate}

% ===== 5. ELABORACIÓN DE RESULTADOS ==================================
\section{ELABORACIÓN DE RESULTADOS}\label{sec:resultados}

\subsection{Identificación de las dimensiones medidas}

La Figura~\ref{fig:croquis} reproduce el croquis de la pieza con la nomenclatura
utilizada. Las letras A a J corresponden a los lados del contorno (vista
frontal), K, L y M a los diámetros de los agujeros, y N y O al espesor y a la
altura total leídos en la vista lateral.

\begin{figure}[H]
	\centering
	\begin{tikzpicture}[scale=0.14,
			lbl/.style={circle,fill=darkblue,text=white,inner
					sep=1.3pt,font=\sffamily\bfseries\scriptsize},
			hlbl/.style={font=\sffamily\bfseries\small,text=darkblue},
			tit/.style={font=\sffamily\small\itshape,text=mediumblue}]
		% --- Vista frontal
		\filldraw[fill=lightblue,draw=darkblue,line width=0.9pt,line join=round]
		(0,0) -- (50,0) -- (50,20) -- (40,20) -- (40,30) -- (50,30) -- (50,40) --
		(15,40) -- (15,16) -- (0,16) -- cycle;
		\foreach \x/\y/\n in {7.5/8/K, 37/8/M, 26/30/L}{
				\filldraw[fill=white,draw=darkblue,line width=0.8pt] (\x,\y) circle
				(5);
				\draw[mediumblue,dashed] (\x-5,\y) -- (\x+5,\y);
				\node[hlbl] at (\x,\y+2.6) {\n};
			}
		\node[lbl] at (25,-3.8) {A};
		\node[lbl] at (-3.8,8) {B};
		\node[lbl] at (7.5,19.6) {C};
		\node[lbl] at (11.2,28) {D};
		\node[lbl] at (32.5,43.8) {E};
		\node[lbl] at (53.8,35) {F};
		\node[lbl] at (45,33.4) {G};
		\node[lbl] at (36.8,25) {H};
		\node[lbl] at (45,16.6) {I};
		\node[lbl] at (53.8,10) {J};
		\node[tit] at (25,-9.5) {Vista frontal};
		% --- Vista lateral
		\filldraw[fill=lightblue,draw=darkblue,line width=0.9pt] (66,0) rectangle
		(77.65,40);
		\node[lbl] at (81.6,20) {O};
		\node[lbl] at (71.8,-3.8) {N};
		\node[tit] at (71.8,-9.5) {Vista lateral};
	\end{tikzpicture}
	\caption{Croquis de la pieza con la nomenclatura empleada (esquema aproximado,
		no a escala exacta).}\label{fig:croquis}
\end{figure}

\subsection{Registro de mediciones}

Las lecturas obtenidas con los tres instrumentos se presentan en la
Tabla~\ref{tab:datos}. Las lecturas en pulgadas se anotaron como fracciones de
$1/128$ in y se incluye su equivalente decimal.

\begin{table}[!htb]
	\centering
	\caption{Mediciones de la pieza con regla, vernier y
		micrómetro}\label{tab:datos}
	\small
	\setlength{\tabcolsep}{4pt}
	\rowcolors{2}{white}{lightblue}
	\begin{tabular}{c l r r c r r}
		\rowcolor{darkblue}
		\encab{Lado}    & \encab{Descripción (según croquis)} & \encab{Regla}   &
		\encab{Vernier} & \encab{Vernier}                     & \encab{Vernier} &
		\encab{Micrómetro}
		\\
		\rowcolor{darkblue}
		                &                                     & \encab{(mm)}    &
		\encab{(mm)}    & \encab{(x/128 in)}                  & \encab{(in)}    &
		\encab{(mm)}
		\\
		A               & Base inferior                       & 49              &
		50,00
		                & 252/128                             & 1,9688          & --
		\\
		B               & Lado izquierdo inferior             & 16              &
		16,30
		                & 82/128                              & 0,6406          &
		16,18
		\\
		C               & Escalón horizontal                  & 16              &
		14,90
		                & 75/128                              & 0,5859          & --
		\\
		D               & Lado izquierdo superior             & 24              &
		23,30
		                & 118/128                             & 0,9219          & --
		\\
		E               & Lado superior                       & 33              &
		34,05
		                & 172/128                             & 1,3438          & --
		\\
		F               & Lado derecho superior               & 10              &
		10,00
		                & 51/128                              & 0,3984          & --
		\\
		G               & Ranura, lado superior               & 10              &
		10,00
		                & 51/128                              & 0,3984          & --
		\\
		H               & Ranura, fondo                       & 10              &
		10,00
		                & 51/128                              & 0,3984          & --
		\\
		I               & Ranura, lado inferior               & 10              & 9,90
		                & 50/128                              & 0,3906          & --
		\\
		J               & Lado derecho inferior               & 20              &
		20,05
		                & 101/128                             & 0,7891          & --
		\\
		K               & Diámetro de agujero                 & 10              &
		10,00           & 51/128                              & 0,3984          & --
		\\
		L               & Diámetro de agujero                 & 10              &
		10,10
		                & 51/128                              & 0,3984          & --
		\\
		M               & Diámetro de agujero                 & 10              &
		8,60            & 44/128                              & 0,3438          & --
		\\
		N               & Espesor (vista lateral)             & 11              &
		11,40
		                & 58/128                              & 0,4531          &
		11,65
		\\
		O               & Altura total (vista lateral)        & 40              &
		40,05           & 202/128                             & 1,5781          & --
		\\
	\end{tabular}
\end{table}

\subsection{Análisis de resultados}

\subsubsection{Comparación entre la regla y el vernier}

Dado que el vernier tiene una resolución veinte veces mejor que la regla
(\SI{0.05}{\milli\metre} frente a \SI{1}{\milli\metre}), se tomó su lectura como
referencia de trabajo para las quince dimensiones y se calculó la diferencia
$\Delta = x_{\text{regla}} - x_{\text{vernier}}$ y el error relativo
$|\Delta|/x_{\text{vernier}}$ (Tabla~\ref{tab:comparacion}).

\begin{table}[!htb]
	\centering
	\caption{Diferencia entre las lecturas de regla y
		vernier}\label{tab:comparacion}
	\small
	\rowcolors{2}{white}{lightblue}
	\begin{tabular}{c r r r r r}
		\rowcolor{darkblue}
		\encab{Lado}            & \encab{Regla (mm)}      & \encab{Vernier (mm)} &
		\encab{$\Delta$
			(mm)}
		                        &
		\encab{$|\Delta|$ (mm)} & \encab{Error rel. (\%)}
		\\
		A                       & 49                      & 50,00                &
		$-1{,}00$
		                        & 1,00                    & 2,00
		\\
		B                       & 16                      & 16,30                &
		$-0{,}30$
		                        & 0,30                    & 1,84
		\\
		C                       & 16                      & 14,90                &
		$+1{,}10$
		                        & 1,10                    & 7,38
		\\
		D                       & 24                      & 23,30                &
		$+0{,}70$
		                        & 0,70                    & 3,00
		\\
		E                       & 33                      & 34,05                &
		$-1{,}05$               & 1,05                    & 3,08
		\\
		F                       & 10                      & 10,00                &
		$0{,}00$
		                        & 0,00                    & 0,00
		\\
		G                       & 10                      & 10,00                &
		$0{,}00$
		                        & 0,00                    & 0,00
		\\
		H                       & 10                      & 10,00                &
		$0{,}00$                & 0,00                    & 0,00
		\\
		I                       & 10                      & 9,90                 &
		$+0{,}10$
		                        & 0,10                    & 1,01
		\\
		J                       & 20                      & 20,05                &
		$-0{,}05$
		                        & 0,05                    & 0,25
		\\
		K                       & 10                      & 10,00                &
		$0{,}00$
		                        & 0,00                    & 0,00
		\\
		L                       & 10                      & 10,10                &
		$-0{,}10$
		                        & 0,10                    & 0,99
		\\
		M                       & 10                      & 8,60                 &
		$+1{,}40$
		                        & 1,40                    & 16,28
		\\
		N                       & 11                      & 11,40                &
		$-0{,}40$               & 0,40                    & 3,51
		\\
		O                       & 40                      & 40,05                &
		$-0{,}05$
		                        & 0,05                    & 0,12
		\\
	\end{tabular}
\end{table}

\begin{figure}[H]
	\centering
	\begin{tikzpicture}
		\begin{axis}[
				width=\linewidth, height=6.4cm,
				symbolic x coords={A,B,C,D,E,F,G,H,I,J,K,L,M,N,O},
				xtick=data, enlarge x limits=0.05,
				ymin=-1.6, ymax=1.8, ytick={-1.5,-1,-0.5,0,0.5,1,1.5},
				ylabel={$\Delta = x_{\mathrm{regla}} - x_{\mathrm{vernier}}$ (mm)},
				xlabel={Lado / diámetro},
				bar width=7pt,
				ymajorgrids, grid style={borderblue!70},
				label style={font=\sffamily\small}, tick label
				style={font=\sffamily\small},
				axis x line*=bottom, axis y line*=left,
				extra y ticks=0, extra y tick labels={}, extra y tick
				style={grid=major,grid style={solid,black!60}}]
			\addplot[ybar,fill=mediumblue,draw=darkblue] coordinates {
					(A,-1.00) (B,-0.30) (C,1.10) (D,0.70) (E,-1.05) (F,0) (G,0) (H,0)
					(I,0.10) (J,-0.05) (K,0) (L,-0.10) (M,1.40) (N,-0.40) (O,-0.05)};
			\addplot[sharp plot,no marks,dashed,thick,red!70!black] coordinates
				{(A,0.5) (O,0.5)};
			\addplot[sharp plot,no marks,dashed,thick,red!70!black] coordinates
				{(A,-0.5) (O,-0.5)};
		\end{axis}
	\end{tikzpicture}
	\caption{Diferencia entre las lecturas de regla y vernier por dimensión. Las
		líneas discontinuas marcan $\pm\SI{0.5}{\milli\metre}$, la incertidumbre de
		lectura esperada para una regla graduada en
		milímetros.}\label{fig:diferencias}
\end{figure}

El resumen estadístico de las diferencias se presenta en la
Tabla~\ref{tab:estadistica}.

\begin{table}[!htb]
	\centering
	\caption{Resumen estadístico de $\Delta$ (regla $-$ vernier),
		$n=15$}\label{tab:estadistica}
	\small
	\rowcolors{2}{white}{lightblue}
	\begin{tabular}{l r}
		\rowcolor{darkblue}
		\encab{Indicador}                                  & \encab{Valor}        \\
		Diferencia media con signo                         & $+0{,}02$ mm         \\
		Diferencia absoluta media $\overline{|\Delta|}$    & 0,42 mm              \\
		Diferencia absoluta media sin el diámetro M        & 0,35 mm              \\
		Desviación estándar de $\Delta$                    & 0,66 mm              \\
		Diferencia absoluta máxima                         & 1,40 mm (diámetro M) \\
		Mediciones con $|\Delta|\le\SI{0.5}{\milli\metre}$ & 10 de 15 (67\,\%)    \\
		Error relativo medio                               & 2,6\,\%              \\
	\end{tabular}
\end{table}

De estos resultados se desprenden varias observaciones:

\begin{itemize}
	\item \textbf{No hay un sesgo sistemático.} La diferencia media con signo es
	      prácticamente nula ($+\SI{0.02}{\milli\metre}$) porque la regla lee unas
	      veces
	      por encima y otras por debajo del vernier. Esto sugiere que las
	      discrepancias
	      provienen de errores aleatorios de lectura (paralaje, alineación del
	      cero,
	      estimación a simple vista) más que de un defecto de la escuadra.
	\item \textbf{Las dimensiones cortas coinciden bien.} Los lados F, G, H y el
	      diámetro K dieron lecturas idénticas, y I, J, L y O difieren en
	      $\SI{0.10}{\milli\metre}$ o menos. Cuando la dimensión coincide con un
	      número
	      entero de milímetros, la regla y el vernier coinciden; la regla pierde
	      exactitud
	      cuando el valor cae entre dos divisiones.
	\item \textbf{Las mayores desviaciones aparecen en A, C, D, E y M.} Cinco de
	      las quince dimensiones superan la incertidumbre esperada de
	      $\pm\SI{0.5}{\milli\metre}$ (Figura~\ref{fig:diferencias}). Las
	      diferencias de C
	      ($+1{,}10$) y E ($-1{,}05$) son casi iguales y de signo opuesto porque
	      son
	      tramos consecutivos de un mismo lado: un error al ubicar con la regla el
	      vértice
	      del escalón se traslada con signo contrario a ambos tramos, mientras que
	      su suma
	      casi coincide (49,00 mm con la regla y 48,95 mm con el vernier). Algo
	      similar
	      ocurre con B y D. Es un error típico al medir con regla tramos
	      consecutivos: se
	      compensa en la suma, pero no en cada tramo.
	\item \textbf{El diámetro M es atípico.} El vernier dio
	      $\SI{8.60}{\milli\metre}$ frente a unos $\SI{10}{\milli\metre}$ de los
	      otros dos agujeros. La causa más probable es una mala posición de las
	      mordazas interiores (que dan una cuerda menor que el diámetro si no se
	      colocan en el diámetro real) o algún borde con rebaba. Se recomienda
	      repetir esta medición antes de dar el valor como definitivo.
\end{itemize}

\subsubsection{Comparación con el micrómetro}

Solo dos dimensiones, B (lado) y N (espesor), pudieron medirse con el
micrómetro. Por su resolución de \SI{0.01}{\milli\metre} y su rango de
\SI{25}{\milli\metre}, se toma como valor de referencia (Tabla~\ref{tab:micro}).

\begin{table}[!htb]
	\centering
	\caption{Errores de regla y vernier respecto del micrómetro}\label{tab:micro}
	\small
	\rowcolors{2}{white}{lightblue}
	\begin{tabular}{c r r r r r r r}
		\rowcolor{darkblue}
		\encab{Lado}        & \encab{Micróm.} & \encab{Regla}       & \encab{$E_a$
			                                                              regla} &
		\encab{$E_r$ regla} & \encab{Vernier} & \encab{$E_a$ vern.} & \encab{$E_r$
			                                                              vern.}
		\\
		\rowcolor{darkblue}
		                    & \encab{(mm)}    & \encab{(mm)}        & \encab{(mm)}
		                    &
		\encab{(\%)}        & \encab{(mm)}    & \encab{(mm)}        & \encab{(\%)}
		\\
		B                   & 16,18           & 16                  & $-0{,}18$    &
		1,11
		                    & 16,30           & $+0{,}12$           & 0,74
		\\
		N                   & 11,65           & 11                  & $-0{,}65$    &
		5,58
		                    & 11,40           & $-0{,}25$           & 2,15
		\\
	\end{tabular}
\end{table}

En ambos casos el vernier estuvo más cerca del micrómetro que la regla (errores
absolutos de $0{,}12$ y $\SI{0.25}{\milli\metre}$ frente a $0{,}18$ y
$\SI{0.65}{\milli\metre}$), como cabe esperar de su mejor resolución. Sin
embargo, la diferencia de $\SI{0.25}{\milli\metre}$ en el espesor N equivale a
cinco veces la resolución del vernier, por lo que no puede atribuirse solo al
instrumento. Entre las causas probables están una fuerza de cierre de las
mordazas distinta a la del trinquete del micrómetro, una ligera inclinación de
las mordazas respecto de las caras, la presencia de óxido o rebabas, o
simplemente que las dos mediciones no se hicieron exactamente en el mismo punto
de la placa. Con solo dos puntos no es posible concluir estadísticamente cuál es
el error típico del vernier; solo puede afirmarse que, en este caso, su error
fue mayor que su resolución.

\subsubsection{Cierre geométrico del contorno}

Un modo de verificar la consistencia de un conjunto de mediciones es comprobar
que las dimensiones \emph{cierran} el contorno. Según el croquis, se debe
cumplir que la altura del lado izquierdo ($B+D$) coincida con la del derecho
($F+H+J$) y con la altura total O, y que la base A sea igual a la suma de los
tramos superiores ($C+E$). Los resultados se muestran en la
Tabla~\ref{tab:cierre}.

\begin{table}[!htb]
	\centering
	\caption{Verificación del cierre geométrico (valores en mm)}\label{tab:cierre}
	\small
	\rowcolors{2}{white}{lightblue}
	\begin{tabular}{l r r}
		\rowcolor{darkblue}
		\encab{Comprobación}                        & \encab{Regla} &
		\encab{Vernier}                                                     \\
		Altura izquierda $B+D$                      & 40,00         & 39,60
		\\
		Altura derecha $F+H+J$                      & 40,00         & 40,05
		\\
		Altura total O                              & 40,00         & 40,05
		\\
		\textbf{Diferencia (izquierda $-$ derecha)} & \textbf{0,00} &
		\textbf{$-0{,}45$}                                                  \\
		Tramo superior $C+E$                        & 49,00         & 48,95
		\\
		Base A                                      & 49,00         & 50,00
		\\
		\textbf{Diferencia ($C+E$ $-$ $A$)}         & \textbf{0,00} &
		\textbf{$-1{,}05$}                                                  \\
	\end{tabular}
\end{table}

Con la regla el contorno cierra exactamente, pero eso no debe interpretarse como
mayor exactitud: al redondear cada lado al milímetro más cercano, los errores se
compensan. Con el vernier, en cambio, aparecen descuadres de
$\SI{0.45}{\milli\metre}$ en vertical y de $\SI{1.05}{\milli\metre}$ en
horizontal, mayores que la incertidumbre combinada de sumar tres o dos lecturas
de $\pm\SI{0.05}{\milli\metre}$. Esto indica que alguna lectura del vernier
(posiblemente la de la base A o la del escalón C) estuvo afectada por chaflanes,
rebabas o por un apoyo inadecuado de las mordazas, y refuerza la recomendación
de repetir las mediciones críticas.

% Conclusiones y recomendaciones

% =========================== 6. CUESTIONARIO =========================
\section{CUESTIONARIO}

\subsection*{A. Identifique y explique el uso de 3 instrumentos por cada unidad
	de medida derivada}

\begin{enunciado}
	A.- Identifica y explique su uso de 3 instrumentos por cada unidad de medida
	derivada.
\end{enunciado}

Se seleccionaron seis unidades derivadas de amplio uso en un laboratorio o
taller de manufactura (Tabla~\ref{tab:cuestionarioA}).

	{\small
		\setlength{\tabcolsep}{5pt}
		\begin{longtable}{>{\raggedright\arraybackslash\bfseries}p{2.45cm}
				>{\raggedright\arraybackslash}p{3.3cm}
				>{\raggedright\arraybackslash}p{8.4cm}}
			\caption{Instrumentos por unidad derivada}\label{tab:cuestionarioA}
			\\
			\rowcolor{darkblue}\encab{Unidad derivada}       & \encab{Instrumento}
			                                                 & \encab{Explicación de
				                                                   su uso}
			\\
			\endfirsthead
			\rowcolor{darkblue}\encab{Unidad derivada}       & \encab{Instrumento}
			                                                 &
			\encab{Explicación de su uso}
			\\
			\endhead
			\rowcolor{lightblue}
			Fuerza\newline(newton, N)                        & Dinamómetro
			                                                 & Mide fuerzas por la
			deformación de un resorte calibrado (o de un sensor en los
			modelos digitales). Se usa para conocer fuerzas de tracción o de sujeción,
			por
			ejemplo la necesaria para extraer una pieza de un útil.
			\\
			\rowcolor{lightblue}
			                                                 & Celda de carga
			                                                 & Transductor con galgas
			extensométricas que convierte la deformación de un
			elemento metálico en una señal eléctrica proporcional a la fuerza. Está
			integrada en prensas, balanzas industriales y máquinas de ensayo.
			\\
			\rowcolor{lightblue}
			                                                 & Máquina universal de
			ensayos                                          & Aplica cargas
			controladas de tracción, compresión o flexión y registra fuerza frente a
			deformación, con lo que se determinan resistencia y ductilidad de los
			materiales.
			\\
			Presión\newline(pascal, Pa)                      & Manómetro de Bourdon
			                                                 & Un tubo curvo elástico
			se deforma con la presión interna y mueve una aguja
			sobre una escala. Es el instrumento típico para circuitos neumáticos e
			hidráulicos, compresores y prensas.
			\\
			                                                 & Transductor de presión
			                                                 & Sensor (normalmente
			piezorresistivo) que entrega una
			señal eléctrica proporcional a la presión, apto para registro continuo y
			control
			automático de procesos.
			\\
			                                                 & Vacuómetro
			                                                 & Mide presiones
			inferiores a la atmosférica. Se usa en sistemas
			de sujeción por vacío, en moldeo y en procesos de desgasificado.
			\\
			\rowcolor{lightblue}
			Par\newline(newton metro, N$\cdot$m)             & Llave dinamométrica
			                                                 & Permite apretar
			pernos con un par especificado y avisa (con un clic o una indicación) al
			alcanzarlo. Es imprescindible en ensamble para no sobre-apretar ni dejar
			juntas
			flojas.
			\\
			\rowcolor{lightblue}
			                                                 & Destornillador
			dinamométrico                                    & Versión de menor
			capacidad para
			tornillería pequeña y componentes sensibles, como conjuntos electrónicos o
			mecánicos de precisión.
			\\
			\rowcolor{lightblue}
			                                                 & Sensor (transductor) de
			par                                              & Mide el par transmitido
			por un eje
			mediante galgas o principios magnetoelásticos. Se usa en bancos de prueba
			de
			motores, husillos y herramientas.
			\\
			Velocidad angular\newline(rad/s; en taller, rpm) & Tacómetro óptico o de
			contacto                                         & Indica las revoluciones
			por minuto de un husillo, eje o disco
			abrasivo. Sirve para ajustar la velocidad de corte y verificar que una
			máquina
			gira a la velocidad programada.
			\\
			                                                 & Estroboscopio
			                                                 & Emite destellos a
			frecuencia ajustable; cuando esta coincide con la del giro,
			el objeto parece detenido. Permite medir frecuencia de rotación o de
			vibración
			sin contacto.
			\\
			                                                 & Encoder rotativo
			                                                 & Genera pulsos por
			vuelta y se usa en máquinas CNC para realimentar posición y
			velocidad de ejes y husillos.
			\\
			\rowcolor{lightblue}
			Volumen\newline(metro cúbico, m$^3$)             & Probeta graduada
			                                                 & Recipiente con escala
			para medir volúmenes de líquidos como refrigerantes,
			aceites de corte o soluciones de limpieza.
			\\
			\rowcolor{lightblue}
			                                                 & Pipeta / bureta
			                                                 & Permite dosificar o
			trasvasar volúmenes pequeños con mayor
			exactitud que la probeta, por ejemplo al preparar mezclas o reactivos de
			laboratorio.
			\\
			\rowcolor{lightblue}
			                                                 & Matraz aforado
			                                                 & Recipiente calibrado
			para un único volumen, empleado al preparar soluciones de
			concentración conocida.
			\\
			Tensión eléctrica\newline(voltio, V)             & Multímetro digital
			                                                 & Mide tensión, y también
			corriente y resistencia. Se usa para diagnosticar
			circuitos de máquinas, sensores y fuentes de alimentación.
			\\
			                                                 & Voltímetro
			                                                 & Instrumento dedicado a
			medir diferencias de potencial; se
			conecta en paralelo con el elemento a medir y es común en tableros de
			control.
			\\
			                                                 & Osciloscopio
			                                                 & Muestra la forma de
			onda de una tensión en el tiempo, lo que permite analizar
			ruido, frecuencia y señales de sensores y variadores.
			\\
		\end{longtable}}

\subsection*{B. Identifique y explique 5 instrumentos de medición usados en un
	laboratorio de manufactura}

\begin{enunciado}
	B.- Identifica y explique 5 instrumentos de medición que se usan en un
	laboratorio de manufactura.
\end{enunciado}

\subsubsection*{1. Pie de rey (vernier)}
Mide longitudes exteriores, interiores y profundidades con resolución típica de
$\SI{0.05}{\milli\metre}$ (o $\SI{0.02}{\milli\metre}$ en algunos modelos). Es
el instrumento de uso más general del taller: sirve para el control dimensional
rápido de piezas mecanizadas, para verificar diámetros y para comprobar
escalones, como se hizo en este laboratorio.

\subsubsection*{2. Micrómetro de exteriores}
Se basa en un tornillo de paso conocido y alcanza resolución de
$\SI{0.01}{\milli\metre}$ (o $\SI{0.001}{\milli\metre}$ en modelos digitales).
Se emplea para medir diámetros de ejes, espesores de placas y otras dimensiones
donde la tolerancia es estrecha. Cada micrómetro cubre un rango de
\SI{25}{\milli\metre} y se elige el que contenga la dimensión, por ejemplo
0--25, 25--50 mm.

\subsubsection*{3. Reloj comparador}
Instrumento de comparación: no mide la dimensión absoluta, sino la
\emph{desviación} respecto de un valor de referencia. Un palpador desplaza un
mecanismo de engranajes o un sensor que mueve la aguja, con resolución típica de
$\SI{0.01}{\milli\metre}$. Montado en un soporte magnético, se usa para
comprobar planitud, paralelismo, excentricidad y alabeo, y para centrar piezas
en un torno o en una fresadora.

\subsubsection*{4. Bloques patrón}
Piezas de acero o cerámica de caras paralelas y espesor calibrado con alta
exactitud, que pueden adherirse entre sí para formar longitudes específicas. No
miden por sí mismos: son \emph{patrones} con los que se calibran y verifican
vernier, micrómetros y comparadores, y con los que se ajusta el cero de un
comparador. Aseguran la trazabilidad de las mediciones del taller.

\subsubsection*{5. Rugosímetro}
Un palpador de punta muy fina recorre la superficie y convierte las
irregularidades en una señal eléctrica, de la que se calculan parámetros como la
rugosidad media $R_a$ (en \si{\micro\metre}). Se usa para verificar que el
acabado superficial logrado por torneado, fresado o rectificado cumple lo
especificado en el plano.

\medskip
Además de estos instrumentos, un laboratorio de manufactura suele contar con
calibres pasa/no pasa, durómetros, proyectores de perfiles y máquinas de
medición por coordenadas (CMM).

\subsection*{C. Mediciones en un proceso de torneado}

\begin{enunciado}
	C.- Cuando usted realice un proceso de torneado, ¿qué mediciones puede
	realizar?
	Explique por qué.
\end{enunciado}

En el torneado la pieza gira y la herramienta avanza, por lo que se miden tanto
las \emph{dimensiones de la pieza} como los \emph{parámetros del proceso}. La
Tabla~\ref{tab:torneado} las resume.

\begin{table}[!htb]
	\centering
	\caption{Mediciones posibles en un proceso de torneado}\label{tab:torneado}
	\small
	\rowcolors{2}{white}{lightblue}
	\begin{tabularx}{\linewidth}{>{\raggedright\arraybackslash\bfseries}p{3.2cm}
			>{\raggedright\arraybackslash}p{3.6cm} X}
		\rowcolor{darkblue}
		\encab{Medición}                                & \encab{Instrumento}
		                                                & \encab{¿Por qué se mide?}
		\\
		Diámetros exteriores                            & Vernier, micrómetro de
		exteriores                                      & Es la dimensión
		principal: se compara con el plano para decidir si se necesita otra pasada o
		si
		la pieza ya está dentro de tolerancia.
		\\
		Diámetros interiores (agujeros)                 & Vernier con mordazas
		interiores,
		micrómetro de interiores, calibre de interiores & Verifica el mandrinado y
		el
		ajuste con el eje que se acoplará.
		\\
		Longitudes y escalones                          & Vernier, profundímetro
		                                                & Controla las cotas
		axiales, hombros y ranuras.
		\\
		Excentricidad y alabeo                          & Reloj comparador con base
		magnética
		                                                & Antes de cortar, se
		comprueba que la pieza esté centrada en el plato; de lo
		contrario el diámetro no será concéntrico.
		\\
		Rugosidad superficial                           & Rugosímetro
		                                                & El acabado depende del
		avance y del radio
		de punta de la herramienta; se
		comprueba que cumpla el plano.
		\\
		Ángulos y conos                                 & Goniómetro, regla de senos
		                                                & Verifica chaflanes y
		conicidad.
		\\
		Roscas                                          & Peine de roscas,
		micrómetro de roscas,
		calibres pasa/no pasa                           & Comprueba paso y diámetro
		para
		asegurar que la rosca ensamble.
		\\
		Velocidad de giro                               & Tacómetro
		                                                & Permite calcular la
		velocidad de corte real y ajustarla al material y a la
		herramienta.
		\\
	\end{tabularx}
\end{table}

Los parámetros del proceso se relacionan con las dimensiones medidas. La
velocidad de corte $v_c$ depende del diámetro $D$ (en mm) y de la velocidad de
giro $n$ (en rpm), y la rugosidad teórica $R_a$ depende del avance $f$ y del
radio de punta $r_\varepsilon$ de la herramienta (ambos en mm):
\[
	v_c = \frac{\pi\,D\,n}{1000}\ \ [\text{m/min}],
	\qquad
	R_a \approx \frac{f^{2}}{32\,r_\varepsilon}\times 1000\ \ [\si{\micro\metre}]
\]
Por ejemplo, para $D=\SI{50}{\milli\metre}$ y $n = 800$ rpm, $v_c = \pi\cdot
	50\cdot 800/1000 \approx \SI{125.7}{\metre\per\minute}$. Otra relación útil es
que, al dar una profundidad de corte $a_p$ (medida en el radio), el diámetro se
reduce en $2\,a_p$, por lo que medir el diámetro tras cada pasada permite
corregir el tambor graduado del carro transversal y compensar el desgaste de la
herramienta.

\begin{nota}[Seguridad]
	Las mediciones de dimensiones con vernier o micrómetro deben hacerse con el
	husillo \textbf{detenido}: medir una pieza en giro es peligroso para el
	operador
	y falsea la lectura. Asimismo, antes de arrancar el torno se debe retirar la
	llave del plato y sujetar el cabello y la ropa suelta.
\end{nota}

% ======================= 7. ENSAYOS VIRTUALES ========================
\section{ENSAYOS VIRTUALES}

Los siguientes simuladores gratuitos permiten practicar la lectura de vernier y
micrómetro en línea y complementan el trabajo realizado en el laboratorio:

\begin{enumerate}[label=\textbf{\arabic*.}]
	\item \textbf{Vernier Calipers Model (versión en mm).} Simulador interactivo
	      del Ministerio de Educación de Singapur / Open Source Physics para
	      practicar la
	      lectura de un vernier en escala métrica moviendo las mordazas. \\
	      \url{https://iwant2study.moe.edu.sg/lookangejss/01_measurement/ejss_model_AAPTVernierCalipermmversion/}
	\item \textbf{Micrometer Model (pantalla completa).} Simulador de micrómetro
	      de exteriores en el que se combina la lectura del manguito y del tambor.
	      \\
	      \url{https://iwant2study.moe.edu.sg/resources/lookangejss/01_measurement/ejss_model_Micrometer02/}
	\item \textbf{Practice on Reading a Vernier Caliper (MiniPhysics).}
	      Explicación paso a paso y entrenador interactivo con lecturas
	      aleatorias. \\
	      \url{https://www.miniphysics.com/practice-on-reading-a-vernier-caliper.html}
	\item \textbf{How to Read a Micrometer Screw Gauge (MiniPhysics).} Guía de
	      lectura del micrómetro con entrenador interactivo. \\
	      \url{https://www.miniphysics.com/how-to-read-a-micrometer-screw-gauge.html}
\end{enumerate}

% ==================== 8. ACCIDENTES RELACIONADOS =====================
\section{ACCIDENTES RELACIONADOS}

Los errores de medición y de conversión de unidades, así como los descuidos en
el manejo de máquinas de taller, han producido accidentes graves. Se sugieren
los siguientes casos y documentos:

\begin{enumerate}[label=\textbf{\arabic*.}]
	\item \textbf{Mars Climate Orbiter (1999).} La sonda de la NASA se perdió al
	      llegar a Marte porque un archivo de software de tierra entregaba datos
	      de los
	      propulsores en unidades inglesas (libras-fuerza) cuando el sistema de
	      navegación
	      esperaba unidades métricas (newtons). Es un ejemplo de cómo una
	      conversión mal
	      gestionada puede echar a perder una misión entera. \\
	      Comunicado de la NASA:
	      \url{https://nssdc.gsfc.nasa.gov/planetary/text/mco_pr_19991110.txt} \\
	      Resumen de \emph{Physics Today}:
	      \url{https://physicstoday.aip.org/news/mars-climate-orbiter-lost} \\
	      Informe de la junta de investigación (Fase I):
	      \url{https://www.dcs.gla.ac.uk/~johnson/Mars/MCO_report.pdf}
	\item \textbf{«Gimli Glider», vuelo 143 de Air Canada (1983).} Un Boeing 767
	      se quedó sin combustible en pleno vuelo porque la carga se calculó con
	      un factor
	      de conversión incorrecto (libras en lugar de kilogramos). La tripulación
	      logró
	      un aterrizaje de emergencia sin víctimas mortales. \\
	      \url{https://en.wikipedia.org/wiki/Gimli_Glider}
	\item \textbf{Seguridad en el uso de tornos y máquinas herramienta.} Guía de
	      la Universidad de Cornell sobre los riesgos de atrapamiento,
	      protecciones y
	      precauciones en el torneado, relevantes cuando se miden piezas cerca de
	      una
	      máquina. \\
	      \url{https://ehs.cornell.edu/campus-health-safety/occupational-safety/tool-and-machine-safety/metalworking-lathe-safety}
\end{enumerate}

\textbf{Riesgos propios del laboratorio de metrología.} Aunque en esta práctica
no hay maquinaria en movimiento, existen riesgos menores: cortes por bordes o
rebabas de la pieza, caída de la pieza o de los instrumentos sobre los pies o
las manos, y daños en los instrumentos por golpes, suciedad u oxidación, que
alteran su calibración. La medición de una pieza en una máquina siempre debe
hacerse con la máquina detenida.

% ========================= 9. BIBLIOGRAFÍA ===========================
\section{BIBLIOGRAFÍA}

\begin{enumerate}[label={[\arabic*]},ref=\arabic*,leftmargin=2.2em,itemsep=5pt]
	\item\label{b:si} Bureau International des Poids et Mesures. (2019). \emph{Le
		      Système international d'unités / The International System of Units}
	      (9.ª ed.).
	      BIPM. \url{https://www.bipm.org/si-brochure}
	\item\label{b:vim} Joint Committee for Guides in Metrology. (2012).
	      \emph{International vocabulary of metrology — Basic and general concepts
		      and
		      associated terms (VIM)} (JCGM 200:2012, 3.ª ed.). BIPM.
	\item\label{b:ley} Congreso de la República del Perú. (1982). \emph{Ley N.º
		      23560: Sistema Legal de Unidades de Medida del Perú}.
	      \url{http://www.congreso.gob.pe/ntley/Imagenes/Leyes/23560.pdf}
	\item\label{b:inacal} Congreso de la República del Perú. (2014). \emph{Ley N.º
		      30224: Ley que crea el Sistema Nacional para la Calidad y el Instituto
		      Nacional
		      de Calidad}.
	\item Groover, M. P. (2007). \emph{Fundamentos de manufactura moderna:
		      materiales, procesos y sistemas} (3.ª ed.). McGraw-Hill.
	\item Kalpakjian, S., \& Schmid, S. R. (2014). \emph{Manufactura, ingeniería y
		      tecnología} (7.ª ed.). Pearson Educación.
	\item National Aeronautics and Space Administration. (1999). \emph{Mars
		      Climate Orbiter Mishap Investigation Board — Phase I Report} y
	      comunicado de
	      prensa del 10 de noviembre de 1999.
	      \url{https://nssdc.gsfc.nasa.gov/planetary/text/mco_pr_19991110.txt}
	\item Cornell University, Environment, Health and Safety. \emph{Metalworking
		      Lathe Safety Awareness Guide}.
	      \url{https://ehs.cornell.edu/campus-health-safety/occupational-safety/tool-and-machine-safety/metalworking-lathe-safety}
	\item Open Source Physics @ Singapore (Hwang, F.-K., Wee, L. K.).
	      \emph{Vernier Calipers Model} y \emph{Micrometer Model}.
	      \url{https://iwant2study.moe.edu.sg/}
\end{enumerate}

\end{document}
